% Section: Introduction — P8 commit 2.
% Numerical claims verified against:
%   - results/p7_8_solver_h1_n24/h1_analysis_combined_n24.json
%   - results/p5_h1_verdict/h1_analysis.json
%   - results/p7_8_kdv_gate/gate_result.json
%   - ODE_PDE_PRE_REG.md (pre-registration)
%   - PRE_REG_AMENDMENT.md (A1--A11)
% All non-numerical claims are cited from refs in paper/references.bib.
% Methodology detail (BAS commitments, lineage, amendments, deferred
% Kuramoto/KdV extensions) lives in Methods (sec:methods); the intro
% states motivation, gap, what-we-did, the headline, and contributions.

\section{Introduction}
\label{sec:intro}

The question of whether variational quantum machine-learning (QML)
models offer a useful advantage over classical neural networks
remains contested~\cite{schuld2019quantum,havlicek2019supervised,
caro2022generalization}. Schuld and Killoran argue that the field
should pursue \emph{regime-dependent} advantages rather than blanket
quantum-supremacy claims~\cite{schuld2022right}: a quantum model is
``useful'' if it solves a specific class of problems better than the
best matched classical model at the same parameter budget.
Physics-informed neural networks (PINNs)~\cite{raissi2019physics,
karniadakis2021physics,lu2021deepxde} are a natural substrate for this
question, because the differential-equation residual loss is a single
scalar that depends on the model class through a clean
derivative-of-parameterized-function interface and admits
capacity-matched classical baselines. The emerging family of
\emph{quantum PINNs}~\cite{perezsalinas2020dataru,kyriienko2021solving,
markidis2022physics,berger2025teqpinn,sedykh2024quantum,
farea2025qcpinn,ahmed2024rfqrc} carries a theoretically motivated
regime-dependence claim: the data-reuploading Fourier series accessible
to a quantum circuit~\cite{schuld2021effect} should grant an inductive
bias for smooth, periodic dynamics. Yet the empirical picture is
fragmentary --- most quantum-PINN papers report a single task, often
without matched classical baselines or pre-registered acceptance
criteria --- leaving the central question \emph{when, and on what kinds
of problems, does a quantum PINN outperform a capacity-matched
classical PINN?} unanswered at any rigorous level.

Bowles, Ahmed, and Schuld recently catalogued the methodological gaps
that make QML benchmarks unreliable~\cite{bowles2024better,
huang2021power} --- undertuned classical baselines, mismatched
parameter counts, best-of-many-seeds reporting, and isolated wins in
place of statistically defensible contrasts --- and propose a set of
commitments (pre-registration, matched budgets, paired-bootstrap
inference, an underfit guard, and a known-structure skyline) that any
advantage claim should satisfy. To our knowledge, no published
quantum-PINN benchmark has implemented this framework on a hardness
ladder spanning both ordinary and partial differential equations. This
paper does so.

Following a pre-registration committed before any results were
generated (Methods~\ref{sec:methods}; Appendix~A), we test the
hypothesis that quantum PINNs achieve a regime-dependent advantage on
smooth/periodic versus broadband/multiscale dynamical systems. Four
quantum-PINN families --- data-reuploading~\cite{perezsalinas2020dataru},
the trainable-embedding QPINN with classical-FNN and with QNN
encoders~\cite{berger2025teqpinn}, and the hybrid
QCPINN~\cite{farea2025qcpinn} --- are run against capacity-matched
classical PINN and non-liquid Neural-ODE baselines on four ODEs
(\lv{}, \vdp{}, Lorenz, \fhn{}) and four PDEs (heat, smooth Burgers,
\ac{}, shock Burgers) at three seeds each (24 cells total), with every
system pre-registered as smooth/periodic or broadband/multiscale.

\textbf{The pre-registered hypothesis is FALSIFIED.} The solver-task
PRIMARY verdict at $n=24$ is $\ddiff = \dsmooth - \dbroad = -0.084$
(paired-bootstrap 95\% CI $[-0.278, +0.061]$, including zero), and the
forecaster-task PRIMARY verdict at $n=9$ is $\ddiff = -0.501$ (CI
$[-0.804, -0.244]$, excluding zero in the \emph{broadband-favoring}
direction). Three additional sensitivity points --- a default-Adam
combined verdict, a both-sides hyperparameter-optimized verdict, and
the original-pre-registration forecaster verdict
(Section~\ref{sec:verdict}) --- return the same qualitative result. The
point estimate even flips sign as the broadband bin is expanded from
six to twelve cells, hinting (without statistical significance) at the
\emph{inverse} of the original hypothesis. The robust conclusion is the
absence of a statistically distinguishable regime-dependent
quantum-PINN advantage at this matched-capacity, matched-budget scale.

Within the falsification, the expanded ladder surfaces one positive
sub-finding worth reporting on its own terms. On \fhn{}, the
trainable-embedding quantum PINN with QNN encoder
(\texttt{te\_qpinn\_qnn})~\cite{berger2025teqpinn} attains a
seed-to-seed standard deviation roughly twenty-four times tighter and a
seed-mean error roughly two-and-a-half times lower than the matched
classical PINN (Section~\ref{sec:solver}) --- a regime-specific
structural edge on stiff fast--slow dynamics, achieved with zero
classical parameters in the ansatz and at the same training budget as
every other cell.

The paper's contributions are fourfold. \emph{Methodologically}, it
delivers an open-source quantum-PINN benchmarking framework that
implements the full Bowles--Ahmed--Schuld
protocol~\cite{bowles2024better} --- pre-registration with falsifiable
acceptance rules, matched parameter budgets, paired-bootstrap
inference, underfit and skyline guards, and symmetric hyperparameter
sensitivity --- backed by a script that mechanically gates every cited
number against a committed JSON file, with twenty-two pre-registration
amendments documented openly (Appendix~A). \emph{Empirically}, it
falsifies the regime-dependent advantage hypothesis across two PRIMARY
verdicts and three layered sensitivity points, covering all eight
hardness-ladder systems and both pre-registered tasks, while
identifying the \texttt{te\_qpinn\_qnn} structural sub-finding at
matched compute. \emph{Mechanistically}, it cross-tabulates the
per-cell advantage gap against trainability and expressivity scalars
(Kullback--Leibler divergence to a Haar-random ensemble, accessible
Fourier $K_{\max}$~\cite{schuld2021effect}, and gradient variance),
finding a leave-one-out-robust but underpowered trend consistent with
the inverted pattern (Section~\ref{sec:mechanism}). Finally, the full
code, data, and reproduction pipeline are released at
\url{https://github.com/shawngibford/qlnn}, with the reproducibility
commitments following NeurIPS program guidance~\cite{pineau2021improving}.

\subsection*{What this paper does \emph{not} claim}

To be precise about scope: this paper does not claim a quantum
advantage; it claims its \emph{absence} at this matched-capacity,
matched-budget scale. It does not claim that the H1 hypothesis is false
in general --- only that the pre-registered acceptance threshold is not
crossed on this hardness ladder. It does not claim that the
broadband-leaning point estimate sign is reliable; the confidence
interval crosses zero on every solver-side sensitivity point. And it
does not assert generalization from $n=24$ cells to the full QML/PDE
landscape; the follow-up experiments named in
Section~\ref{sec:discussion} are precisely what that would require.

The rest of the paper proceeds as follows: Methods
(Section~\ref{sec:methods}), solver- and forecaster-task results
(Sections~\ref{sec:solver}--\ref{sec:forecaster}), the aggregated H1
verdict and sensitivity analysis (Section~\ref{sec:verdict}), the
mechanism cross-tabulation (Section~\ref{sec:mechanism}), and
discussion and conclusions
(Sections~\ref{sec:discussion}--\ref{sec:conclusions}).
